\section{Centrikus ütközés mechanikai modellje}

\subsection*{Két test centrikus ütközésének modellje}
A mechanikai modell a következő ábrán látható.
\begin{figure}[H]
    \centering
    \resizebox{\textwidth}{!}{
    \begin{tikzpicture}[>=stealth, font=\small]
        \def\ang{20}
        \def\Rone{0.8}
        \def\Rtwo{0.55}
        
        \tikzset{
            ball1/.style={thick, ball color=gray!40},
            ball2/.style={thick, ball color=green!60!black!60},
            vel/.style={->, thick, blue!80!black},
            com/.style={yellow, draw=black, fill=yellow},
            c1cent/.style={cyan, draw=black, fill=cyan},
            c2cent/.style={orange, draw=black, fill=orange}
        }

        % --- Ütközés előtt ---
        \begin{scope}[shift={(0,0)}]
            \draw[->] (\ang-180:1.6) -- (\ang:2.8) node[right] {$\mathbf{n}$};
            
            % ghosts
            \draw[gray, opacity=0.2, fill=gray!10] (-0.8, {-0.8*tan(\ang)-0.3}) circle (\Rone);
            \draw[gray, opacity=0.4, fill=gray!20] (-0.4, {-0.4*tan(\ang)-0.1}) circle (\Rone);

            \coordinate (C1) at (0, 0);
            \coordinate (C2) at (\ang:1.35);
            \coordinate (S) at (\ang:0.85);

            \draw[ball1] (C1) circle (\Rone);
            \draw[ball2] (C2) circle (\Rtwo);
            
            \node at (-0.7, 0.8) {$m_1$};
            \node at (2.05, 0.6) {$m_2$};

            \filldraw[c1cent] (C1) circle (1.5pt);
            \filldraw[c2cent] (C2) circle (1.5pt);
            \filldraw[com] (S) circle (1.5pt);

            \draw[vel] (C1) -- +(60:2) node[left] {$\mathbf{c}_{S1}$};
            \draw[vel] (C2) -- +(-30:1.2) node[below] {$\mathbf{c}_{S2}$};
            \draw[vel] (S) -- +(45:1.5) node[above] {$\mathbf{c}_S$};

            \node at (0.75, -2) {Ütközés előtt};
        \end{scope}

        % --- Ütközés közben ---
        \begin{scope}[shift={(5.5,0)}]
            \draw[->] (\ang-180:1.6) -- (\ang:2.8) node[right] {$\mathbf{n}$};
            
            \coordinate (C1) at (0, 0);
            \coordinate (C2) at (\ang:1.15);
            \coordinate (S) at (\ang:0.72);

            \draw[ball1] (C1) circle (\Rone);
            \draw[ball2] (C2) circle (\Rtwo);

            % Collision waves
            \draw[gray, thick] (0.55, 0.4) to[bend left=20] (0.7, 0.05);
            \draw[gray, thick] (0.45, 0.45) to[bend left=20] (0.6, -0.1);
            \draw[gray, thick] (0.35, 0.5) to[bend left=20] (0.5, -0.2);

            \filldraw[c1cent] (C1) circle (1.5pt);
            \filldraw[c2cent] (C2) circle (1.5pt);
            \filldraw[com] (S) circle (1.5pt);

            % Vector C1
            \coordinate (V1) at ([shift={(65:2.2)}]C1);
            \draw[vel] (C1) -- (V1);
            \coordinate (NEND) at (\ang:2.8);
            \coordinate (P1) at ($(C1)!(V1)!(NEND)$);
            \draw[dashed, gray!80] (V1) -- (P1);

            % Vector S
            \coordinate (VS) at ([shift={(45:1.5)}]S);
            \draw[vel] (S) -- (VS) node[above] {$\mathbf{c}_S$};
            \coordinate (PS) at ($(S)!(VS)!(NEND)$);
            \draw[dashed, gray!80] (VS) -- (PS);

            % Vector C2
            \coordinate (V2) at ([shift={(-10:1.3)}]C2);
            \draw[vel] (C2) -- (V2);
            \coordinate (P2) at ($(C2)!(V2)!(NEND)$);
            \draw[dashed, gray!80] (V2) -- (P2);

            \node at (0.6, -2) {Ütközés közben};
        \end{scope}

        % --- Ütközés után ---
        \begin{scope}[shift={(11,0)}]
            \draw[->] (\ang-180:1.6) -- (\ang:2.8) node[right] {$\mathbf{n}$};
            
            \coordinate (C1) at (0, 0);
            \coordinate (C2) at (\ang:1.35);
            \coordinate (S) at (\ang:0.85);

            \draw[ball1] (C1) circle (\Rone);
            \draw[ball2] (C2) circle (\Rtwo);

            \filldraw[c1cent] (C1) circle (1.5pt);
            \filldraw[c2cent] (C2) circle (1.5pt);
            \filldraw[com] (S) circle (1.5pt);

            \draw[vel] (C1) -- +(80:1.8) node[left] {$\mathbf{v}_{S1}$};
            \draw[vel] (C2) -- +(5:2.2) node[below] {$\mathbf{v}_{S2}$};
            \draw[vel] (S) -- +(45:1.5) node[above] {$\mathbf{c}_S$};

            \node at (0.75, -2) {Ütközés után};
        \end{scope}
    \end{tikzpicture}
    }
    \caption{Két test centrikus ütközésének részletes folyamatábrája.}
\end{figure}
A jelölt mennyiségek:
\begin{itemize}
    \item $m_1, m_2$: A két test tömege.
    \item $\mathbf{n}$: Az ütközési egyenes iránya (normális).
    \item $\mathbf{c}_{S1}, \mathbf{c}_{S2}$: Az ütközés előtti sebességvektorok.
    \item $\mathbf{v}_{S1}, \mathbf{v}_{S2}$: Az ütközés utáni sebességvektorok.
    \item $\mathbf{c}_S$: A közös tömegközéppont (súlypont) állandó sebessége.
\end{itemize}

\subsection*{A newtoni modell alapfeltevései}
A centrikus ütközés newtoni modelljének fontosabb alapfeltevései:
\begin{itemize}
    \item \textbf{Ütközés időbeli lefutása:} Az ütközés lezajlási ideje (a kontaktus ideje) $\Delta t \to 0$, azaz az ütközés pillanatszerű.
    \item \textbf{Sebességkomponensek megváltozása:} Mivel az erők (belső ütközési erők) impulzusa nagyon véges rövid idő alatt történik, csak az ütközési egyenes menti sebességkomponensek szenvednek véges ugrásszerű megváltozást. Az ütközési egyenesre merőleges (érintőirányú) sebességkomponensek nem változnak meg (súrlódásmentes ütközés).
    \item A külső erők impulzusa elhanyagolható a belső (ütközési) erők impulzusához képest.
\end{itemize}

\newpage
\subsection*{Impulzusok és a közös súlypont sebessége}
Az ütközés során a testekből álló zárt rendszer összimpulzusa megmarad, mivel a külső erők hatása elhanyagolható ezen a rövid $\Delta t$ idő alatt. Vagyis az ütközés előtti ($p_1 + p_2$) és utáni ($\tilde{p}_1 + \tilde{p}_2$) impulzusok összege egyenlő:  $m_1 v_1 + m_2 v_2 = m_1 \tilde{v}_1 + m_2 \tilde{v}_2$. 

Ennek következménye az is, hogy a két testből álló rendszer közös súlypontjának sebessége ($v_S$) az ütközés során \textbf{állandó marad}. Így mind a kompressziós, mind a restitúciós szakasz alatt ez a sebesség változatlan.
